\documentclass[11pt,a4paper]{article}
\usepackage[utf-8]{inputenc}
\usepackage[margin=1in]{geometry}
\usepackage{amsmath}
\usepackage{graphicx}
\usepackage{xcolor}
\usepackage{listings}
\usepackage{hyperref}
\usepackage{fancyhdr}
\usepackage{float}

\lstset{
    language=bash,
    basicstyle=\ttfamily\small,
    breaklines=true,
    frame=single,
    backgroundcolor=\color{gray!10},
    keywordstyle=\color{blue},
    commentstyle=\color{gray},
}

\pagestyle{fancy}
\fancyhf{}
\rhead{\thepage}
\lhead{NMPDE Projects: CI/CD and Testing Infrastructure}

\title{\textbf{NMPDE Projects: Comprehensive Testing and CI/CD Documentation}}
\author{Advanced Methods for Scientific Computing}
\date{February 2026}

\begin{document}

\maketitle

\section*{Abstract}
This document provides a comprehensive overview of the testing and continuous integration infrastructure developed for the NMPDE space-time adaptivity project. It covers the automated verification pipeline, scalability testing framework, model validation mechanisms, and Azure cloud deployment architecture. The documentation serves as a reference guide for developers and stakeholders involved in the project.

\newpage
\tableofcontents

\newpage

\section{Introduction and Project Overview}

\subsection{Project Scope}
The NMPDE (Numerical Methods for Partial Differential Equations) projects repository implements advanced numerical algorithms for solving partial differential equations, with a specific focus on space-time adaptivity techniques. The primary deliverable is the `space-time-adaptivity` executable, built using the \textit{deal.II} finite element library and implemented in C++.

\subsection{Key Components}
The project structure includes:
\begin{itemize}
    \item \textbf{Source Code} (\texttt{src/}): Core implementation in C++, including the Heat equation solver and manufactured solution tests
    \item \textbf{Mesh Files} (\texttt{mesh/}): Gmsh-generated mesh files for various refinement levels
    \item \textbf{Build System}: CMake-based configuration using deal.II's integration
    \item \textbf{Testing Infrastructure} (\texttt{scripts/}): Bash and Python utilities for verification, validation, and performance analysis
    \item \textbf{CI/CD Pipelines} (\texttt{.github/workflows/}): GitHub Actions workflows for continuous integration
    \item \textbf{Infrastructure as Code} (\texttt{terraform/}): Terraform configurations for Azure cloud provisioning
\end{itemize}

\subsection{Technology Stack}
\begin{itemize}
    \item \textbf{Language}: C++17 with deal.II library
    \item \textbf{Build System}: CMake (\geq 3.12)
    \item \textbf{Container Platform}: Docker (quay.io/pjbaioni/amsc\_mk:2025)
    \item \textbf{CI/CD}: GitHub Actions
    \item \textbf{Cloud Infrastructure}: Microsoft Azure with Terraform
    \item \textbf{Testing Framework}: CTest (integrated with CMake)
    \item \textbf{Validation}: Python 3 with pyvista for VTU/PVTu mesh analysis
\end{itemize}

\newpage

\section{Testing Infrastructure Overview}

\subsection{Testing Philosophy}
The project employs a multi-tiered testing approach:
\begin{enumerate}
    \item \textbf{Unit Tests}: CTest-based tests defined in \texttt{CMakeLists.txt}
    \item \textbf{Verification Tests}: Automated validation of compilation, execution, and correctness
    \item \textbf{Model Validation Tests}: Computation-intensive tests verifying physical conservation laws
    \item \textbf{Scalability Tests}: Performance benchmarking across varying thread counts
    \item \textbf{Artifact Generation}: Automated production of visualization-ready output files (VTU/PVTu format)
\end{enumerate}

\subsection{Build Configuration}
The project uses a standard CMake workflow:

\begin{lstlisting}
cmake_minimum_required(VERSION 3.12.0)
project(space-time-adaptivity LANGUAGES CXX C)
include(common/cmake-common.cmake)

add_executable(space-time-adaptivity src/space-time-adaptivity.cpp src/Heat.cpp)
deal_ii_setup_target(space-time-adaptivity)

enable_testing()
add_test(NAME space-time-adaptivity COMMAND space-time-adaptivity)
add_test(NAME model_validation COMMAND space-time-adaptivity)
\end{lstlisting}

Key CMake options:
\begin{itemize}
    \item \texttt{CMAKE\_BUILD\_TYPE=Release}: Optimized compilation for performance
    \item \texttt{DDEAL\_II\_ALLOW\_TESTS=ON}: Enable deal.II test support
\end{itemize}

\subsection{Local Testing Execution}

\subsubsection{Verification Testing}
The verification script configures, builds, and runs all registered tests:

\begin{lstlisting}
#!/bin/bash
./scripts/run_verification.sh [build_dir]
\end{lstlisting}

This script:
\begin{enumerate}
    \item Configures the project with Release optimization
    \item Builds all targets in parallel
    \item Executes CTest with parallel execution on $(nproc)-1 cores
    \item Reports test results with full failure diagnostics
\end{enumerate}

\subsubsection{Model Validation Testing}
The model validation harness provides sophisticated verification capabilities:

\begin{lstlisting}
TEST_CMD='./build/bin/project5_generate_results > model_results.csv' \
  RESULTS_FILE=model_results.csv TOL=1e-6 \
  PY_VALIDATOR=scripts/validate_model.py \
  ./scripts/run_model_validation.sh build
\end{lstlisting}

Environment variables for customization:
\begin{itemize}
    \item \texttt{VENV\_DIR}: Python virtual environment location (default: \texttt{venv})
    \item \texttt{BUILD\_DIR}: Build directory containing executables
    \item \texttt{RESULTS\_FILE}: Output file containing numerical results
    \item \texttt{TEST\_CMD}: Custom command to generate results
    \item \texttt{PY\_VALIDATOR}: Path to Python validation script
    \item \texttt{TOL}: Relative error tolerance for validation (default: \texttt{1e-6})
\end{itemize}

\subsubsection{Scalability Testing}
Thread-based scalability analysis with timing measurements:

\begin{lstlisting}
TEST_EXEC=build/space-time-adaptivity ./scripts/run_scalability.sh build
\end{lstlisting}

Output format (CSV):
\begin{lstlisting}
threads,seconds,cmd
1,45.32,build/space-time-adaptivity --threads=1
2,23.18,build/space-time-adaptivity --threads=2
4,12.45,build/space-time-adaptivity --threads=4
8,7.33,build/space-time-adaptivity --threads=8
16,4.92,build/space-time-adaptivity --threads=16
\end{lstlisting}

\newpage

\section{GitHub Actions CI/CD Pipeline}

\subsection{Architecture Overview}
The GitHub Actions infrastructure consists of two coordinated workflows:
\begin{itemize}
    \item \textbf{Verification CI}: Automatic on every push and pull request
    \item \textbf{Scalability Tests}: Manual dispatch workflow with optional executable parameter
\end{itemize}

Both workflows execute within the standardized Docker container \texttt{quay.io/pjbaioni/amsc\_mk:2025}, ensuring reproducibility across deployment environments.

\subsection{Verification Workflow}

The verification workflow (\texttt{.github/workflows/verification.yml}) implements a two-stage pipeline:

\subsubsection{Stage 1: Build and Test}
\begin{itemize}
    \item \textbf{Environment}: Ubuntu latest with containerized deal.II dependencies
    \item \textbf{Modules}: Loads \texttt{gcc-glibc} and \texttt{dealii} via environment modules
    \item \textbf{Build}: Parallel CMake configuration and compilation
    \item \textbf{Testing}: Executes CTest with parallel execution strategy
    \item \textbf{Artifacts}: Uploads entire build directory for downstream model validation
\end{itemize}

\subsubsection{Stage 2: Model Validation}
\begin{itemize}
    \item \textbf{Dependency}: Requires successful build-and-test completion
    \item \textbf{Environment}: Isolated job with download of build artifacts
    \item \textbf{Execution}: Runs \texttt{scripts/run\_model\_validation.sh}
    \item \textbf{Validation}: Executes Python-based model correctness verification
\end{itemize}

Key Features:
\begin{itemize}
    \item Automatic triggering on push and pull requests
    \item Parallel build optimization ($(nproc)-1 threads)
    \item Artifact retention for cross-job communication
    \item Comprehensive failure diagnostics with full output
\end{itemize}

\subsection{Scalability Workflow}

The scalability workflow (\texttt{.github/workflows/scalability.yml}) provides configuration flexibility:

\begin{itemize}
    \item \textbf{Trigger}: Manual dispatch (\texttt{workflow\_dispatch})
    \item \textbf{Input Parameter}: \texttt{test\_executable} (default: \texttt{build/space-time-adaptivity})
    \item \textbf{Build Step}: Clean rebuild of project in Release mode
    \item \textbf{Execution}: Runs \texttt{scripts/run\_scalability.sh} with provided executable
    \item \textbf{Output}: Generates \texttt{scalability\_results.csv}
\end{itemize}

Workflow dispatch allows flexibility such as:
\begin{lstlisting}
# Run with custom benchmark executable
TEST_EXEC=build/custom-benchmark ./scripts/run_scalability.sh build

# Or default to ctest scalability tests
./scripts/run_scalability.sh build
\end{lstlisting}

\newpage

\section{Model Validation Framework}

\subsection{Validation Architecture}
The model validation system provides sophisticated verification for numerical solutions through multiple analysis pathways:

\subsubsection{Input Format Support}
\begin{enumerate}
    \item \textbf{CSV Numeric Files}: Direct numerical result files
    \item \textbf{VTU Files}: Individual time-step visualization files from deal.II's DataOut
    \item \textbf{PVTu Files}: Master parallel visualization files (typically used for analysis)
\end{enumerate}

\subsubsection{Analysis Methodology}
For VTU/PVTu format outputs, the validator:
\begin{enumerate}
    \item Reads mesh data using PyVista library
    \item Extracts the \texttt{solution} field from point or cell data
    \item Computes volume-weighted mean (mass integral)
    \item Compares against known reference value
    \item Reports relative error and pass/fail status
\end{enumerate}

\subsection{Validation Script Implementation}

The Python validator (\texttt{scripts/validate\_model.py}) provides multiple validation strategies:

\begin{lstlisting}[language=python]
def compute_mean_from_vtu(path: Path):
    """Computes the mean of the 'solution' data from VTU/PVTu."""
    import pyvista as pv
    mesh = pv.read(str(path))
    
    if 'solution' in mesh.point_data:
        return float(np.mean(mesh.point_data['solution']))
    if 'solution' in mesh.cell_data:
        return float(np.mean(mesh.cell_data['solution']))
    
    # Fallback to first available data array
    if mesh.point_data:
        key = list(mesh.point_data.keys())[0]
        return float(mesh.point_data[key][0])
    raise ValueError('No suitable data arrays found')
\end{lstlisting}

\subsection{Physical Conservation Properties}
The default validation target assumes the \texttt{FunctionU0} initial condition over the unit domain $[0,1]^3$. The expected volume-weighted mean is:

\[
\text{Expected Mass} = \int_{\Omega} u_0(x) \, dx = \frac{1}{216}
\]

The numerical scheme with homogeneous Neumann boundary conditions should conserve this quantity. Validation succeeds when:

\[
\left|\frac{\text{Computed Mean} - \frac{1}{216}}{\frac{1}{216}}\right| < \text{TOL}
\]

where TOL is configurable (default: $10^{-6}$).

\newpage

\section{Azure Cloud Infrastructure and Deployment}

\subsection{Infrastructure Overview}
The Azure infrastructure provides scalable cloud resources for distributed testing and performance monitoring. Terraform manages all provisioned resources, ensuring reproducibility and version control.

\subsection{Resource Architecture}

\subsubsection{Resource Group}
\textbf{Name}: \texttt{nmpde-projects-rg}
\begin{itemize}
    \item Central container for all resources
    \item Ensures logical grouping and lifecycle management
    \item Supports role-based access control (RBAC)
\end{itemize}

\subsubsection{Container Registry (ACR)}
\textbf{Name}: \texttt{nmpdeprojectsacr}
\begin{itemize}
    \item Tier: Basic (suitable for development/testing)
    \item Admin access enabled for GitHub Actions authentication
    \item Mirrors container images (e.g., quay.io/pjbaioni/amsc\_mk:2025)
    \item Integration point between GitHub Actions and Azure deployment
\end{itemize}

\subsubsection{Storage Account}
\textbf{Name}: \texttt{nmpdeprojectsstorage}
\begin{itemize}
    \item Tier: Standard with Local Redundancy (LRS)
    \item Blob containers for organizational separation:
    \begin{enumerate}
        \item \texttt{performance-reports}: JSON and CSV test outputs
        \item \texttt{build-artifacts}: Compiled binaries and libraries
        \item \texttt{test-results}: CTest and benchmark results
    \end{enumerate}
    \item Access: Private (requires authentication)
\end{itemize}

\subsubsection{Log Analytics Workspace}
\textbf{Configuration}: PerGB2018 pricing tier
\begin{itemize}
    \item Centralizes monitoring and diagnostic data
    \item Enables performance analysis and troubleshooting
    \item Integrates with Azure Monitor for alerting
\end{itemize}

\subsection{Infrastructure as Code (Terraform)}

\subsubsection{Provider Configuration}
\begin{lstlisting}
terraform {
  required_version = ">= 1.0"
  required_providers {
    azurerm = {
      source  = "hashicorp/azurerm"
      version = "~> 3.0"
    }
  }
}
provider "azurerm" {
  features {}
  subscription_id = var.subscription_id
}
\end{lstlisting}

\subsubsection{Deployment Process}
\begin{enumerate}
    \item \textbf{Initialization}: \texttt{terraform init} (downloads providers)
    \item \textbf{Planning}: \texttt{terraform plan} (previews changes)
    \item \textbf{Application}: \texttt{terraform apply} (provisions resources)
    \item \textbf{Cleanup}: \texttt{terraform destroy} (removes resources)
\end{enumerate}

\subsubsection{Variable Configuration}
Key Terraform variables:
\begin{itemize}
    \item \texttt{subscription\_id}: Azure subscription identifier
    \item \texttt{resource\_group\_name}: Logical container name
    \item \texttt{location}: Azure region (e.g., eastus)
    \item \texttt{acr\_name}: Container registry identifier
    \item \texttt{storage\_account\_name}: Storage account identifier
\end{itemize}

\subsection{GitHub Actions Integration}

\subsubsection{Service Principal Setup}
Azure resources require GitHub Actions to authenticate. Setup process:

\begin{lstlisting}
# Create service principal with Contributor role
az ad sp create-for-rbac \
  --name "github-actions-nmpde" \
  --role contributor \
  --scopes /subscriptions/5d7fbbd4-db18-4472-915d-de671cc0f22a \
  --sdk-auth
\end{lstlisting}

Output JSON contains credentials stored as GitHub secret \texttt{AZURE\_CREDENTIALS}.

\subsubsection{Secrets Configuration}
Repository secrets required:
\begin{itemize}
    \item \texttt{AZURE\_CREDENTIALS}: Service principal authentication JSON
    \item \texttt{AZURE\_SUBSCRIPTION\_ID}: Subscription identifier
    \item \texttt{AZURE\_RESOURCE\_GROUP}: Resource group name
\end{itemize}

\subsection{Performance Monitoring and Reporting}

The performance testing script (\texttt{scripts/run\_performance\_tests.sh}) collects comprehensive system metrics:

\subsubsection{Collected Metrics}
\begin{itemize}
    \item \textbf{CPU Information}: Core count, architecture, flags (from \texttt{lscpu})
    \item \textbf{Memory}: Total, available, and used capacity (from \texttt{free})
    \item \textbf{Disk}: Storage capacity and utilization (from \texttt{df})
    \item \textbf{Process Limits}: Maximum processes and memory constraints (from \texttt{ulimit})
    \item \textbf{Timing Data}: Execution duration and throughput metrics
\end{itemize}

\subsubsection{Output Format}
Reports are generated in two formats:
\begin{enumerate}
    \item \textbf{JSON Report} (\texttt{performance\_YYYYMMDD\_HHMMSS.json}): Structured data for automated processing
    \item \textbf{Text Report} (\texttt{system\_metrics\_YYYYMMDD\_HHMMSS.txt}): Human-readable system information
\end{enumerate}

Reports are timestamped to enable historical trend analysis and performance regression detection.

\end{document}
